\documentclass[11pt]{article}
\usepackage[a4paper,margin=29mm]{geometry}
\usepackage{amsmath,amssymb,amsthm}
\usepackage{graphicx}
\usepackage{xcolor}
\usepackage{microtype}
\usepackage[hidelinks]{hyperref}
\usepackage{enumitem}

\definecolor{deepblue}{HTML}{173B57}
\definecolor{softgray}{HTML}{F3F5F7}
\newcommand{\Mc}{\mathcal M}
\newcommand{\Pc}{\mathcal P}
\newcommand{\Tc}{\mathcal T_c}
\newtheorem{theorem}{Theorem}

\title{\vspace{-1.5em}\textbf{Vague Mass Escape Destroys Lower Semicontinuity}\\
\large A counterexample for generalized transport costs on locally finite measures}
\author{FA/Banach discovery run \texttt{fa\_banach\_001}\\
\texttt{agent\_lane\_02} \quad Model: \texttt{GPT5.6}}
\date{11 August 2026}

\begin{document}
\maketitle

\begin{center}
\colorbox{softgray}{\parbox{0.91\linewidth}{
\textbf{Status: counterexample likely valid.}
The question in Gozlan--Herry--Peccati has a negative answer as stated, even
for a bounded continuous ordinary transport cost that is linear in the
conditional law.}}
\end{center}

\begin{abstract}
Gozlan, Herry, and Peccati ask whether a lower-semicontinuous generalized
transport cost remains lower semicontinuous after replacing finite measures
with Radon measures finite on balls and weak convergence with vague
convergence.  We give a two-Dirac counterexample on $\mathbb R$.  The mechanism
is mass escape: vague convergence does not preserve total mass, while their
extended transport functional is infinite for measures of unequal mass.
\end{abstract}

\section{Source question}

The source is N. Gozlan, R. Herry, and G. Peccati, \emph{Transport
inequalities for random point measures}, arXiv:2002.04923, later published in
the \emph{Journal of Functional Analysis} \textbf{281} (2021), no.~9,
109141~\cite{GHP}.  On page 10, at the end of Section 2.2.3, the authors ask:

\begin{quote}
``When $c$ is lower semi-continuous, we do not know if
$\mathcal T_c:\Mc_0(E)\times\Mc_0(E)\to[0,\infty]$ is lower
semi-continuous.''
\end{quote}

\begin{figure}[ht]
\centering
\fbox{\includegraphics[width=0.96\linewidth]{figures/open_problem_crop.png}}
\caption{The question in the source paper, page 10.  The preceding paragraph
proves lower semicontinuity for finite measures with the weak topology.}
\end{figure}

\section{Definitions and convention}

Let $E$ be a Polish metric space.  The source denotes by $\Mc_0(E)$ the Radon
measures finite on balls and equips this space with the vague topology,
generated by maps
\[
  \nu\longmapsto\int_E f\,d\nu
\]
for continuous $f$ vanishing outside a ball.  Given
$c:E\times\Pc(E)\to[0,\infty]$, it defines
\[
 \Tc(\nu_2\mid\nu_1)
 =\inf_N\int_E c(x,p_x)\,\nu_1(dx),
 \qquad N(dx,dy)=p_x(dy)\nu_1(dx),
\]
where $N$ ranges over couplings of $\nu_1$ and $\nu_2$.  In particular, the
source explicitly adopts
\begin{equation}\label{eq:mass-convention}
 \Tc(\nu_2\mid\nu_1)=+\infty
 \quad\text{whenever}\quad \nu_1(E)\ne\nu_2(E).
\end{equation}

\section{Proof intuition}

Weak convergence of finite measures sees the constant function $1$, so total
mass cannot disappear.  Vague convergence does not: a unit Dirac mass may run
to infinity and converge vaguely to zero.  Equation~\eqref{eq:mass-convention}
turns precisely this harmless-looking mass loss into an infinite jump of the
transport cost.  Truncating the Euclidean distance makes every prelimit cost
uniformly bounded, while the limiting pair has unequal mass.

\section{Counterexample}

\begin{theorem}\label{thm:counterexample}
There exist a Polish space $E$ and a bounded, jointly continuous function
$c:E\times\Pc(E)\to[0,\infty)$, linear in its second argument, for which
\[
 \Tc:\Mc_0(E)\times\Mc_0(E)\longrightarrow[0,\infty]
\]
is not lower semicontinuous in the product vague topology.
\end{theorem}

\begin{proof}
Take $E=\mathbb R$ with its usual metric, set
\[
  \rho(x,y)=1\wedge|x-y|,
  \qquad
  c(x,p)=\int_{\mathbb R}\rho(x,y)\,p(dy),
\]
and define, for integers $n\ge1$,
\[
  \nu_{1,n}=\delta_0,
  \qquad
  \nu_{2,n}=\delta_n.
\]

First, $c$ satisfies all the advertised regularity.  It is bounded by $1$ and
linear, hence convex, in $p$.  If $x_j\to x$ and $p_j\to p$ weakly in
$\Pc(\mathbb R)$, then
\begin{align*}
 |c(x_j,p_j)-c(x,p)|
 &\le \sup_y|\rho(x_j,y)-\rho(x,y)|
   +\left|\int\rho(x,y)\,(p_j-p)(dy)\right|\\
 &\le |x_j-x|+o(1),
\end{align*}
because $y\mapsto\rho(x,y)$ is bounded and continuous.  Thus $c$ is jointly
continuous.

The unique coupling of $\delta_0$ and $\delta_n$ is
$\delta_{(0,n)}$.  Therefore
\begin{equation}\label{eq:prelimit}
 \Tc(\nu_{2,n}\mid\nu_{1,n})
 =c(0,\delta_n)=1
 \qquad(n\ge1).
\end{equation}

On the other hand, if $f$ is continuous and vanishes outside a ball, then
$f(n)=0$ for every sufficiently large $n$.  Hence
\[
 \nu_{2,n}=\delta_n\longrightarrow0,
 \qquad
 \nu_{1,n}=\delta_0\longrightarrow\delta_0
\]
vaguely in $\Mc_0(\mathbb R)$.  The two limiting measures have total masses
$0$ and $1$, respectively.  By~\eqref{eq:mass-convention},
\begin{equation}\label{eq:limit}
 \Tc(0\mid\delta_0)=+\infty.
\end{equation}
Combining~\eqref{eq:prelimit} and~\eqref{eq:limit} gives
\[
 \Tc(0\mid\delta_0)=+\infty
 >1
 =\liminf_{n\to\infty}\Tc(\delta_n\mid\delta_0),
\]
which is the failure of lower semicontinuity.
\end{proof}

\section{Verification notes}

\begin{enumerate}[leftmargin=2em]
\item Every measure used is Radon and finite on balls.
\item The example uses an ordinary linear transport cost, not a pathological
generalized cost; $c$ is bounded and jointly continuous.
\item Both measures have mass $1$ at every prelimit index, so their coupling
set is nonempty and the cost calculation is exact.
\item Only one of the two unit masses escapes every bounded set.  This makes
the limiting masses unequal and invokes exactly the source's $+\infty$
convention.
\item The example does not contradict the source's preceding theorem for
finite measures with the weak topology.  Weak convergence tests $1$ and thus
preserves total mass; vague convergence is the essential change.
\end{enumerate}

No computation or probabilistic approximation is used.  The packet's separate
verification report repeats the hypothesis audit in checklist form.

\section{Limitations and novelty check}

This counterexample answers the question negatively in its stated generality.
It does not rule out lower semicontinuity on subclasses where vague convergence
is supplemented by tightness or preservation of total mass, nor does it
address alternative unbalanced transport functionals that penalize rather
than forbid mass loss.

A bounded search was performed on 11 August 2026.  We searched the run's
registry, solution, attempt, and proof-gap indexes; the exact arXiv identifier,
paper title, and question phrase; and combinations of ``generalized transport
cost'', ``vague topology'', ``locally finite measures'', ``mass escape'', and
``lower semicontinuity'', including arXiv-focused searches.  The search found
the original arXiv/JFA paper and general work on vague convergence, but no
later paper explicitly answering this exact question.  Novelty is plausible
but not certified by an exhaustive bibliographic review.

\enlargethispage{5\baselineskip}
\section{Human-review recommendation}

\textbf{Recommend promotion as a full negative answer.}  The proof is short
and exact.  A reviewer should principally confirm the source convention
\eqref{eq:mass-convention} and the use of the product vague topology; both are
stated in Section 2.2 of the source.

{\small
\begin{thebibliography}{9}
\bibitem{GHP}
N. Gozlan, R. Herry, and G. Peccati,
\emph{Transport inequalities for random point measures},
Journal of Functional Analysis \textbf{281} (2021), no.~9, 109141;
arXiv:\href{https://arxiv.org/abs/2002.04923}{2002.04923}.
\end{thebibliography}
}

\end{document}
